\subsection{从李代数态射到李群态射的过渡}

引理 1. 设 \(\mathbf{G}\) 是李群芽，\(\mathfrak{h}\) 是 \(\mathbf{L}(\mathbf{G})\) 中具有拓扑补空间的李子代数。\(\mathfrak{gh}\)（分别为 \(\mathfrak{hg}\)）对 \(g\in\mathbf{G}\) 的并集是 \(\mathbf{T}(\mathbf{G})\) 的可积向量子丛。

考虑 \(\mathrm{T}(\mathrm{G})\) 的左平凡化（§ 2，第3号），立即看出，\(g\mathfrak{h}\) 对于 \(g\in \mathrm{G}\) 是向量子丛 \(\mathrm{E}\)（它是 \(\mathrm{T}(\mathrm{G})\) 的子丛）的纤维。设 \(g\in \mathrm{G}\)。\((\mathrm{L}_a)_g\)（其中 \(a\in \mathfrak{h}\)）的集合等于 \(g\mathfrak{h}\)。现在，若 \(a\) 和 \(b\) 属于 \(\mathfrak{h}\)，则 \([\mathrm{L}_a,\mathrm{L}_b] = \mathrm{L}_{[a,b)}\) 且 \([a,b]\in \mathfrak{h}\)。因此 \(\mathrm{E}\) 可积（《可微与解析流形》，R，9.3.3 (iv)）。对于 \(\mathfrak{h}g\)，论证类似。

\(\mathfrak{gh}\)（分别为 \(\mathfrak{hg}\)）的并集的积分叶状结构（《可微与解析流形》，R，9.3.2）称为与 \(\mathfrak{h}\) 相关的 G 的左（分别为右）叶状结构。

定理 1. 设 \(\mathbf{G}\) 和 \(\mathbf{H}\) 是李群芽，\(f\) 是从 \(\mathbf{L}(\mathbf{G})\) 到 \(\mathbf{L}(\mathbf{H})\) 的连续态射。

\begin{enumerate}
\def\labelenumi{(\roman{enumi})}
\item
  存在一个开李子群芽 \(\mathbf{G}'\)，它是 \(\mathbf{G}\) 的开李子群芽，以及一个态射 \(\phi\)，其定义域为 \(\mathbf{G}'\)、陪域为 \(\mathbf{H}\)，使得 \(f = \mathbf{L}(\phi)\)。
\item
  设 \(\mathrm{G}_1, \mathrm{G}_2\) 是 G 的开李子群芽，\(\phi_i\) 是从 \(\mathrm{G}_i\) 到 H 的态射，满足 \(f = \mathrm{L}(\phi_i)\)，其中 \(i = 1, 2\)。则 \(\phi_1\) 与 \(\phi_2\) 在 \(e\) 的某个邻域上重合。设 \(p_1: \mathrm{G} \times \mathrm{H} \to \mathrm{G}\)、\(p_2: \mathrm{G} \times \mathrm{H} \to \mathrm{H}\) 是典范投影。对所有 \((g,h) \in \mathrm{G} \times \mathrm{H}\)，设 \(f_{g,h}\) 是映射 \(ga \mapsto hf(a)\)，从 \(\mathrm{T_g(G)} = g\mathrm{L(G)}\) 到 \(\mathrm{T_h(H)} = h\mathrm{L(H)}\)。考虑 \(\mathrm{T(G)}\) 和 \(\mathrm{T(H)}\) 的左平凡化，立即看出 \(f_{g,h}\) 定义了从 \(p_1^*\mathrm{T(G)}\) 到 \(p_{2}^{*}\mathbf{T}(\mathrm{H})\) 的态射。设 \(a\) 是 \(f\) 的图；它是 \(\mathbf{L}(\mathrm{G}) \times \mathbf{L}(\mathrm{H})\) 的闭李子代数，并以 \(\{0\} \times \mathbf{L}(\mathrm{H})\) 为拓扑补空间。对所有 \((g,h) \in \mathrm{G} \times \mathrm{H}\)，\(f_{g,h}\) 的图是 \((g,h)\) . \(a\)。这些图的并集是 \(\mathrm{T}(\mathrm{G} \times \mathrm{H})\) 的可积向量子丛（引理1）。于是存在（《可微与解析流形》，R，9.3.7）一个开邻域 \(\mathrm{U}\)，包含 \(e_{\mathrm{G}}\) 且位于 \(\mathrm{G}\) 中，以及一个解析映射 \(\phi\)，其定义域为 \(\mathrm{U}\)、陪域为 \(\mathrm{H}\)，使得 \(\phi(e_{\mathrm{G}}) = e_{\mathrm{H}}\)，并且 \(\mathrm{T}_{g}(\phi) = f_{g,\phi(g)}\) 对所有 \(g \in \mathrm{U}\)。特别地，\(\mathrm{T}_{e_{\mathrm{G}}}(\phi) = f\)。
\end{enumerate}

设 \(\mathbf{V}\) 是包含 \(e_{\mathrm{G}}\) 且位于 \(\mathbf{G}\) 中的开邻域，使得对 \((s,t)\in \mathbf{V}\times \mathbf{V}\)，乘积 \(st\) 和 \(\phi (s)\phi (t)\) 有定义且 \(st\in \mathbf{U}\)。考虑映射 \(\alpha_{1},\alpha_{2}\)，它们定义在 \(\mathbf{V}\times \mathbf{V}\) 上、取值于 \(\mathbf{H}\)，并由下式给出：

\[
\alpha_ {1} (s, t) = \phi (t s), \quad \alpha_ {2} (s, t) = \phi (t) \phi (s).
\]

于是 \(\alpha_{1}(t,e) = \phi (t) = \alpha_{2}(t,e)\)。另一方面，固定 V 中的 \(t\)，设 \(\beta_{i}\) 是从 V 到 H 的映射 \(s\mapsto \alpha_i(s,t)\)。则对所有 \(s\in \mathrm{V}\) 及 \(a\in \mathrm{L}(\mathrm{G})\)，

\[
\begin{array}{r l} \mathrm{T} _ {s} (\beta_ {1}) (s a) & = \mathrm{T} _ {t s} (\phi) (t s a) = f _ {t s, \phi (t s)} (t s a) \\ & = \phi (t s) f (a) = f _ {s, \beta_ {1} (s)} (s a) \\ \mathrm{T} _ {s} (\beta_ {2}) (s a) & = \phi (t) \mathrm{T} _ {s} (\phi) (s a) = \phi (t) f _ {s, \phi (s)} (s a) \\ & = \phi (t) \phi (s) f (a) = f _ {s, \beta_ {2} (s)} (s a). \end{array}
\]

因此，根据《可微与解析流形》R，9.3.7，\(\alpha_{1}\) 与 \(\alpha_{2}\) 在 \((e_{\mathrm{G}}, e_{\mathrm{G}})\) 的某个邻域上重合。因此，将 \(\phi\) 限制在 \(e_{\mathrm{G}}\) 的充分小的对称开邻域上，得到李群芽的态射，从而得到 (i)。

设 \(\mathrm{G}_1, \mathrm{G}_2, \phi_1, \phi_2\) 如 (ii) 中，并证明 \(\phi_1, \phi_2\) 在 \(e_{\mathrm{G}}\) 的某个邻域上重合。存在 \(e_{\mathrm{G}}\) 的一个开邻域 W，使得 \(\phi_1(t s) = \phi_1(t)\phi_1(s)\)、\(\phi_2(t s) = \phi_2(t)\phi_2(s)\) 对其中所有 \(s\)、\(t\) 成立。于是，若 \(s \in W\) 且 \(a \in L(G)\)，

\[
\mathrm{T} _ {s} \left(\phi_ {i}\right) (s a) = \phi_ {i} (s) \mathrm{T} _ {e} \left(\phi_ {i}\right) (a) = \phi_ {i} (s) f (a) = f _ {s, \phi _ {i} (s)} (s a)
\]

对 \(i = 1,2\) 均成立。由于 \(\phi_1(e_{\mathrm{G}}) = e_{\mathrm{H}} = \phi_2(e_{\mathrm{G}})\)，由《可微与解析流形》R，9.3.7，\(\phi_1\) 与 \(\phi_2\) 在 \(e_{\mathrm{G}}\) 的某个邻域上重合。

推论 1. 设 G 和 H 是两个李群芽。若 L(G) 与 L(H) 同构，则 G 与 H 局部同构。

这由定理1及 § 1 第10号命题21推出。

推论 2. 设 \(\mathbf{G}\) 是李群芽。若 \(\mathrm{L}(\mathbf{G})\) 是交换的，则 \(\mathbf{G}\) 与加法李群 \(\mathrm{L}(\mathbf{G})\) 局部同构。

加法群 \(\mathbf{L}(\mathbf{G})\) 的李代数与 \(\mathbf{L}(\mathbf{G})\) 同构。因此只需应用推论1。

推论 3. 设 G 是李群。若 L(G) 是交换的，则 G 含有一个交换开子群。

存在 G 的一个交换开李子群芽 U（推论2）。设 V 是 e 的一个邻域，满足 \(V^{2} \subset U\)。则对 V 中的所有 x、y，都有 xy = yx。因此，由 V 生成的 G 的子群是交换的；它显然是开的。

